\section{Security}\label{sec:txid-security}

\subsection{Proof of \texorpdfstring{\Cref{thm:txid-bb}}{the black box
construction theorem}}\label{sec:txid-bb-sec}

We show that the view of a real world adversary $\adv$ during an execution of
$\Pi$ is indistinguishable from the view of an ideal world simulator $\simu$
that interacts with system participants through ideal functionality
$\idfu{\algo{trace}}$. We define a sequence of indistinguishable games that
start with Game 0 representing the real world, and end with Game 7 representing
the ideal world.

We rely on the following security properties of the underlying primitives:

\begin{enumerate}

  \item[(P1)] The registration/credential signature scheme is EUF-CMA secure: no
PPT adversary can produce a valid signature without querying the registration
authority.

  \item[(P2)] The commitment scheme $\ComCommit$ is hiding: commitments are
computationally indistinguishable from random for any PPT adversary. It is also
binding, which the extraction step of (P4) relies on.

  \item[(P3)] $\FE$ and $\FT$ are pseudorandom functions, and remain so against
a distinguisher holding the user's public tracing key $\PTK$. This is what
\Cref{sec:txid-construction} asks for. $\PTK$ is a public parameter of $\FT$
rather than key material, so nothing it carries need be kept secret, and because
it is fixed for the lifetime of the user it is kept out of the output of $\FE$
on purpose. In the instantiation the two statements are
\Cref{lem:txid-fe-prf,lem:txid-ft-prf} respectively.

  \item[(P4)] The proofs and signatures of knowledge for relations $\rel_4$
through $\rel_7$ are simulation extractable in the random oracle model: a valid
witness can be extracted from any accepted proof, and simulated proofs are
indistinguishable from real ones. Extraction proceeds by rewinding via the
forking lemma, which is why the security model of~\Cref{sec:txid-definitions} is
the standalone one rather than universal composability. The discussion there
sets out how straight-line extraction would be obtained instead.

  \item[(P5)] Honest users and honest auditors communicate over a secure
(confidential) channel.

\end{enumerate}

\begin{proof}

  $\simu$ maintains three tables, analogous to $\reg$, $\perm$, and $\txns$ in
$\idfu{\algo{trace}}$:

  \begin{enumerate}

    \item $\regsimu[\user] = (\sk, \utk, \signature)$, logging registered users.

    \item $\permsimu[(\auditor,\user,\epoch)] = \ek$, logging tracing
permissions and epoch keys.

    \item $\txnssimu[(\user,\epoch)]$, logging submitted transactions.

  \end{enumerate}

  \paragraph{Game 0.} The real execution of protocol $\Pi$.

  \paragraph{Game 1.} As Game 0, except $\simu$ simulates all honest
participants (the ledger operator, registration authority, honest users, and
honest auditors) by following $\Pi$'s instructions. Game 1 is indistinguishable
from Game 0 by construction.

  \paragraph{Game 2.} As Game 1, except during registration $\simu$ simulates
the registration authority's output directly from $\idfu{\algo{trace}}$'s
$\IDregister$ calls, rather than executing $\Pi$. For each honest user's
$\IDregister$ request, $\simu$ selects random $(\sk, \utk)$, simulates an
execution of $\Pi$, and stores $\regsimu[\user] = (\sk, \utk, \signature)$.

  For a valid registration request from a corrupt user, $\simu$ issues
$\IDregister$ to $\idfu{\algo{trace}}$ on the user's behalf and returns its
output.  If $\idfu{\algo{trace}}$ registers the user, $\simu$ extracts $(\sk,
\utk)$ using the extractability of the proof (P4), computes $\signature$
following $\Pi$, sends it to $\user$, and records $\regsimu[\user] = (\sk, \utk,
\signature)$.

  The two games diverge only if $\idfu{\algo{trace}}$'s decision on some
$\IDregister$ call differs from $\Pi$'s. The correctness of $\Pi$ rules out
$\idfu{\algo{trace}}$ accepting where $\Pi$ rejects, and the registration
authority's tracking of past requests rules out the reverse, since only a double
registration could cause it, and that never succeeds. So Game 2 is
indistinguishable from Game 1.

  \paragraph{Game 3.} As Game 2, except $\simu$ simulates honest users'
responses to tracing authorisation requests using only $\idfu{\algo{trace}}$'s
output on $\IDgrant$ calls.

  For an $\IDgrant$ call between honest user $\user$ and honest auditor
$\auditor$ in epoch $\epoch$, $\simu$ sets $\permsimu[(\auditor,\user,\epoch)]
\leftarrow \varnothing$ and simulates the transcript with randomly generated
information. By (P5), this is indistinguishable from a real execution.

  For an authorisation request from a corrupt auditor $\auditor$ to honest user
$\user$, $\simu$ forwards it to $\user$, who calls $\IDgrant$. If
$\idfu{\algo{trace}}$ accepts, $\simu$ retrieves $(\sk, \utk, \signature)
\leftarrow \regsimu[\user]$, simulates $\user$'s response following $\Pi$, and
sets $\permsimu[(\auditor,\user,\epoch)] \leftarrow \FE(\utk, \epoch)$.

  This simulated response is indistinguishable from a real one, so Game 3 is
indistinguishable from Game 2.

  \paragraph{Game 4.} As Game 3, except for an authorisation request from a
corrupt auditor for honest user $\user$ and epoch $\epoch$, $\simu$ replaces the
epoch tracing key by a uniform value, and simulates the proof for relation
$\rel_5$.

  Concretely, $\simu$ samples $\ek$ uniformly from the range of $\FE$ and sets
$\permsimu[(\auditor,\user,\epoch)] \leftarrow \ek$. The user's public tracing
key $\PTK_\user$ is sampled once per user, at registration, and is reused in
every grant for that user. Reusing it is not an optimisation but a requirement:
$\PTK$ is certified once and is by construction identical across epochs, so a
simulator that drew a fresh one per grant would be trivially distinguishable by
a corrupt auditor holding grants for two epochs.

  When such a user submits a transaction, $\simu$ computes the tag as
$\FT(\PTK_\user,\ek,\gamma)$ under this substituted key, so that the corrupt
auditor's tracing still succeeds, and simulates the zero-knowledge proof for
relation $\rel_7$. By pseudorandomness of $\FE$ (P3), which holds against a
distinguisher given $\PTK_\user$ and is applied across the epochs for which the
auditor holds no permission, and by (P4), this game is indistinguishable from
Game 3.

  \paragraph{Game 5.} As Game 4, except $\simu$ also simulates transaction
submission using $\idfu{\algo{trace}}$'s output on $\IDsubmit$ requests from
honest users.

  For each honest user's $\IDsubmit$ call in epoch $\epoch$, $\simu$ checks
whether any corrupt auditor can trace the transaction back to its origin. It
does this by issuing $\IDidentify$ queries on behalf of every corrupt auditor
holding a permission for that user, until it finds a match:

  \begin{enumerate}

    \item For each pair of honest user and corrupt auditor $(\user, \auditor)$
with $\permsimu[(\auditor,\user,\epoch)]\neq \bot$, call $\IDidentify$ and
receive $\txns[(\user,\epoch)]$.

    \item If $\txns[(\user,\epoch)] \neq \txnssimu[(\user, \epoch)]$, let
$\algo{TX} = \txns[(\user,\epoch)] \setminus \txnssimu[(\user, \epoch)]$, update
$\txnssimu[(\user, \epoch)] \leftarrow \txns[(\user, \epoch)]$, and output pair
$(\user, \algo{TX})$.

  \end{enumerate}

  Because these two steps run on \emph{every} $\IDsubmit$ call, $\txnssimu$ is
kept in lockstep with $\txns$ for every user that some corrupt auditor can
trace. At most one entry can therefore have grown since the previous call, and
if one has, it belongs to the user who just called $\IDsubmit$ and identifies
them. This leaves exactly two cases. If no entry grew, no corrupt auditor holds
a permission for the calling user this epoch, so $\simu$ never learns who
called. If one grew, say $\txns[(\user^\star,\epoch)] =
\txnssimu[(\user^\star,\epoch)] \cup \{\payload\}$, then a corrupt auditor is
authorised to trace $\user^\star$ and $\simu$ learns both the caller and the
payload. We treat the two cases below.

  \begin{description}

    \item[No corrupt auditor holds a permission for this epoch.] $\simu$ sets
the confidential payload to an arbitrary ciphertext, and instead of computing
the commitment as $\com = \ComCommit(\ck, \utk, \openinfo)$ and the tag as
$\Tag_\gamma = \FT(\PTK,\ek,\gamma)$, samples uniformly at random commitment
$\com^{\star}$ and tag $\Tag_{\gamma}^\star$ and simulates both corresponding
proofs. After Game 4 the epoch keys released to corrupt auditors are uniform and
independent across epochs, so the keys released for preceding and succeeding
epochs carry no information about the key in use here, and $\adv$ therefore
does not hold it. Pseudorandomness of $\FT$ licenses the substitution:
the tags of an epoch are pseudorandom to an adversary holding $\PTK$ but not
$\ek$, which is the position of every corrupt auditor here. So by (P2), (P3),
and (P4), this is indistinguishable from submitting transactions in an execution
of $\Pi$.

    \item[A corrupt auditor holds a valid tracing permission for the honest
user.] $\simu$ obtains pair $(\user, \algo{TX})$, where $\user$ identifies the
honest user and $\algo{TX}$ consists of the payload that $\user$ sent with their
$\IDsubmit$ query. $\simu$ retrieves tuple $(\sk, \utk,\signature) \leftarrow
\regsimu[\user]$ and the substituted epoch key $\ek \leftarrow
\permsimu[(\auditor,\user,\epoch)]$, places that payload in the transaction,
computes the tag as $\FT(\PTK_\user,\ek,\gamma)$ so that the corrupt auditor's
tracing succeeds, and simulates the rest of the transaction as in Game 4.

  \end{description}

  Both cases are indistinguishable from a real execution, so Game 5 is
indistinguishable from Game 4.

  \paragraph{Game 6.} As Game 5, except $\simu$ also simulates transaction
submission for corrupt users, using the output of $\idfu{\algo{trace}}$'s
$\IDsubmit$ interface.

  When a corrupt user $\user$ submits a transaction with a valid tag (a tag with
valid proofs which does not appear in the ledger), $\simu$ uses (P4) to extract
$\user$'s witness, and (P1) to identify $\user$ via $\regsimu$.  $\simu$ then
sends an $\IDsubmit$ request on $\user$'s behalf to $\idfu{\algo{trace}}$ and
returns its output. If $\idfu{\algo{trace}}$ accepts, $\simu$ records the
transaction in $\txns[(\user,\epoch)]$.

  This game is indistinguishable from Game 5 unless $\idfu{\algo{trace}}$
rejects a transaction that $\Pi$ deems valid, which happens only in one of two
cases:

  \begin{description}

    \item[$\user$ is not registered.] This cannot be the case given that each
entry in $\regsimu$ corresponds to a successful $\IDregister$ call to
$\idfu{\algo{trace}}$.

    \item[$\user$ already submitted $\Max$ transactions in the current epoch.]
By (P4) and the determinism of $\FT$, $\Pi$ will not accept a transaction from
$\user$ in this case as the tag will already appear in the ledger.

  \end{description}

  \paragraph{Game 7.} As Game 6, except $\simu$ stops the game if the output of
the $\IDidentify$ interface differs from that of protocol $\Pi$'s tracing
functionality.

  $\simu$ stops if one of two events occurs. Event $e_1$ is a coalition of
corrupt auditors and users tracing a transaction submitted by an honest user
during an epoch even though the user did not grant permissions to any corrupt
auditor for that epoch. Event $e_2$ is an honest auditor failing to trace a
transaction of a corrupt user submitted during an epoch despite holding the
granted authorisation. We show each event is negligible.

  Event $e_1$ does not occur, since the transactions of honest users during an
epoch for which they did not grant tracing permissions to any corrupt auditor
are generated at random.

  Event $e_2$ occurs only with negligible probability, and the argument is the
heart of non-repudiation, so we give it in full. A transaction $\tx$ from a
corrupt user is recorded in $\txnssimu$ and forwarded to $\idfu{\algo{trace}}$
only if it carries a valid tag $\Tag_{\gamma}$, meaning the proof for relation
$\rel_6$ and the proof for relation $\rel_7$ are both valid. By (P4) we extract
a witness from each. From $\rel_6$ we obtain a credential $\signature$ on
$(\G(\sk), \G'(\utk))$ verifying under~\Cref{eq:rel6}, together with an opening
of $\com$ to $\utk$. By (P1) the credential was issued by the registration
authority, so $\utk$ is the tracing key of a registered user, and by the binding
of $\com$ (P2) it is the same $\utk$ that $\rel_7$ speaks about. From $\rel_7$
we obtain $\Tag_{\gamma} = \FT(\PTK, \ek, \gamma)$ with $\PTK = \G'(\utk)$, $\ek
= \FE(\utk, \epoch)$, and
$0 \leq \gamma < \Max$, verifying~\Cref{eq:rel7}. Chaining these, every
transaction the validators accept from a corrupt user carries a tag lying in the
set $\{\FT(\G'(\utk),\FE(\utk,\epoch),\gamma)\}_{\gamma \in [0,\Max)}$ for that
user's own
certified $\utk$, which is exactly the set an authorised auditor enumerates. A
corrupt user therefore cannot place a transaction on the ledger that a permitted
auditor will miss, and in particular the range clause is what rules out the
escape $\gamma \geq \Max$. What remains is that an honest auditor with the right
permission might have accepted a bogus epoch key $\ek^\star$ at grant time,
which by (P4) and (P1) occurs only with negligible probability. So Game 7, which
is the ideal world, is indistinguishable from Game 6. This concludes the proof
of~\Cref{thm:txid-bb}.

\end{proof}

\subsection{Proof of \texorpdfstring{\Cref{lem:txid-fe-prf}}{the
pseudorandomness lemma for the epoch function}}

We show that no polynomial time adversary can distinguish $\epoch \mapsto
\HE(\epoch)^{\utk'}$ from a random function with non-negligible advantage.

\begin{proof}

  Let $\adv$ be an adversary with oracle access to either $\HE(\cdot)^{\utk'}$
or a random function $R: \{0,1\}^* \to \group{1}$, distinguishing the two with
advantage $\varepsilon = \left| \Pr[\adv = 1 \mid \HE(\cdot)^{\utk'}] - \Pr[\adv
= 1 \mid R]\right|$. We use $\adv$ to construct a simulator $\simu$ that breaks
DDH in $\group{1}$.

  $\simu$ receives a DDH challenge $(\gen, \gen^a, \gen^b, Z)$ where $Z =
\gen^{ab}$ (DDH) or $Z \sample \group{1}$ (random). It does the following:

  \begin{enumerate}

    \item Set the secret key implicitly to $\utk'^\star = b$.

    \item For each query $\epoch_i$ to $\HE(\epoch_i)$, draw fresh $r_i, t_i
\sample \Zp$ and set $\HE(\epoch_{i}) = (\gen^a)^{r_i}\cdot\gen^{t_i} =
\gen^{a r_i + t_i}$.

    \item For each query $\epoch_i$ to the PRF oracle, return
$Z^{r_i}\cdot(\gen^{b})^{t_i}$.

  \end{enumerate}

  The pair $(r_i,t_i)$ is stored on the first query on $\epoch_i$ and reused
afterwards, so both oracles are consistent across repeated queries.

  If $Z = \gen^{ab}$, then the response is $\gen^{a b r_i + b t_i} =
\HE(\epoch_i)^{b}$, so $\simu$ has reproduced $\FE$ exactly, with key $b$.

  If $Z = \gen^{z}$ for uniform $z$, then the answer pair for $\epoch_i$ is
$(\gen^{a r_i + t_i}, \gen^{z r_i + b t_i})$. The exponents are the image of the
uniform vector $(r_i,t_i)$ under the matrix $\left(\begin{smallmatrix} a & 1 \\
z & b\end{smallmatrix}\right)$, which is invertible whenever $ab \neq z$. Each
pair is therefore uniform in $\group{1}^2$, and, as the $(r_i,t_i)$ are drawn
independently, the pairs are mutually independent. The view of the distinguisher
is that of a random function, and we lose only the event $ab = z$, of
probability $1/p$.

  Re-randomising with $r_i$ alone would not suffice. It would give answers
$\gen^{z r_i} = \HE(\epoch_i)^{z/a}$, which is again $\FE$, under key $z/a$
rather than $b$, so the two cases would be identically distributed and the
reduction would say nothing. The term $\gen^{t_i}$ is what breaks that
correlation.

  The distinguisher may hold $\PTK$ throughout. Its first component $T =
\gen_2^{\utk'}\cdot\gen_3^{\utk''}$ is a perfectly hiding Pedersen commitment
for $\utk'' \sample \Zp$, so $\simu$ samples it uniformly, and its second
component $\{W_\gamma\}$ is a function of $\utk''$ alone, which $\simu$ draws
itself. Neither touches the challenge. As $\simu$ also plays the registration
authority it chooses the remaining BBS+ generators and the issuing key, and can
produce every other public value of the protocol. This is what lets the lemma be
stated against a distinguisher given the whole of $\PTK$, which
is~\Cref{lem:txid-ft-prf} needs in order to compose with this one.

  An adversary that distinguishes the epoch key from random can then distinguish
between the case where $Z$ is a valid DDH tuple and the case where $Z$ is
random. We then have the following advantage:
%
  \begin{equation*}
%
    \Adv[\adv_{\PRF}(\FE)] \leq \Adv[\adv_\mathsf{DDH}(\simu)] + \frac{1}{p}.
%
  \end{equation*}
%
  The additive term is the event $ab = z$ isolated above, on which the
programmed pairs stop being uniform. Nothing in the bound depends on the number
of random oracle queries, since a fresh pair $(r_i, t_i)$ is drawn for each new
epoch and reused from then on, so no two epochs can collide on a programmed
value
of $\HE$.

\end{proof}

\subsection{Proof of \texorpdfstring{\Cref{lem:txid-ft-prf}}{the
pseudorandomness lemma for the tag function}}

\begin{proof}

  Let $\adv$ distinguish the tag vector of an epoch from random with advantage
$\varepsilon$. We use $\adv$ to build a simulator $\simu$ that breaks
$q$-mDBDHI with $q = \Max$.

  $\simu$ receives a $q$-mDBDHI challenge $(\gen, \ggen_D,
\{\ggen_D^{1/(x+\gamma)}\}_{\gamma\in[0,\Max)},
\{\Delta_{\gamma}\}_{\gamma\in[0,\Max)})$, where either $\Delta_{\gamma} =
\pairing(A, \ggen_D^{1/(x+\gamma)})$ for every $\gamma$ and hidden $A \sample
\group{1}$, or every $\Delta_{\gamma} \sample \group{T}$ independently. It
forwards $\{\ggen_D^{1/(x+\gamma)}\}_{\gamma\in[0,\Max)}$ to $\adv$ as the
public parameters $\{W_\gamma\}$ and $\{\Delta_{\gamma}\}$ as the tag vector,
implicitly setting $\ek^\star = A$ and $\utk''^\star = x$, and echoes $\adv$'s
guess.

  No interpolation over the queried indices is needed, since the challenge
already supplies every $W_\gamma$ directly, matching exactly what an auditor
receives in the real protocol. If the challenge is real, the distribution
$\simu$ hands to $\adv$ is exactly $\{\FT(\PTK,\ek^\star,\gamma)\}_{\gamma\in
[0,\Max)}$ under key $\ek^\star$. If the challenge is random, every entry is
independent and uniform in $\group{T}$. As the domain $[0,\Max)$ is of
polynomial size, indistinguishability of the whole vector from a uniform vector,
which is what the assumption gives, is equivalent to pseudorandomness in the
sense of~\Cref{sec:prelims-prfs}: an adaptive distinguisher can be run on the
vector directly. We then have the following advantage:
%
  \begin{equation*}
%
    \Adv[\adv_{\PRF}(\FT)] \leq
\Adv[\adv_{q\text{-}\mathsf{mDBDHI}}(\simu)].
%
  \end{equation*}
%
  A single instance is enough for the whole lifetime of a user. The epoch keys
of distinct epochs are uniform and mutually independent
by~\Cref{lem:txid-fe-prf}, whose statement hands the distinguisher the whole of
$\PTK$, and it is $\PTK$ alone that the epochs share. The two lemmas therefore
compose without a hybrid over epochs, which is what the $\{W_\gamma\}$ being
public parameters rather than a component of $\ek$ buys us.

  The group elements $\{W_\gamma\}$ do not compromise $\utk''$ itself.
Recovering
$\utk''$ from them would break the $q$-SDH assumption (\Cref{asm:qsdh-bbs}),
which we already assume for the BBS+ signature scheme. And $\utk''$ is fixed
once at registration and never leaves the user, while the $\{W_\gamma\}$ given
to an auditor are certified public parameters, identical in every epoch. So
nothing in a grant links one epoch to another beyond what $\ek$ itself discloses
to auditors who already hold permission for that epoch.

\end{proof}

\subsection{Equivalence of $\rel_7$ and $\rel_7'$}\label{sec:rel7-equiv}

Recall from~\Cref{sec:tag-computation} that $\rel_7$ has instance $\inst_7 =
(\HE, \ggen_D, \gen_2, \gen_3, \gen_4, \com, \Max, \epoch, \Tag_\gamma)$ and
witness $\witness_7 = (\utk', \utk'', \openinfo, \ek, \gamma)$, and holds
exactly when the following does:
%
\begin{equation} \begin{gathered}
%
  \com = \gen_2^{\utk'}\cdot\gen_3^{\utk''}\cdot\gen_4^{\openinfo} \ \wedge \
\ek = \HE(\epoch)^{\utk'} \\ \wedge \ \Tag_{\gamma} = \pairing\big(\ek,\,
\ggen_D^{1/(\utk''+\gamma)}\big) \ \wedge \ 0 \leq \gamma < \Max.
%
\end{gathered} \end{equation}
%
$\rel_7'$ is defined for instance $\inst_7' = (\inst_7, U, \Pi, X, \com_\gamma)$
and witness $\witness_7' = (\witness_7, s, r_\gamma)$, and holds exactly when
the following does:
%
\begin{equation} \begin{gathered}
%
  \com = \gen_2^{\utk'}\cdot\gen_3^{\utk''}\cdot\gen_4^{\openinfo} \ \wedge \ U
= \HE(\epoch)^s \ \wedge \ \Pi = \pairing(U,\ggen_D)^{\utk'} \\ \wedge \ X =
\Tag_\gamma^s \ \wedge \ \Pi = X^{\utk''}\cdot X^{\gamma} \\ \wedge \
\com_\gamma = \gen_3^{\gamma}\cdot\gen_4^{r_\gamma} \ \wedge \ 0 \leq \gamma <
\Max.
%
\end{gathered} \end{equation}
%
For simplicity we argue equivalence for the interactive variant, that is, the
sigma protocol without the Fiat-Shamir transform, and return to the
non-interactive case at the end.

\begin{proof}\leavevmode

  \paragraph{Correctness.} Let $\witness_7$ satisfy $\rel_7$, and sample $s,
r_\gamma \sample \Zp$. Set $U = \HE(\epoch)^{s}$, $\Pi =
\pairing(\ek^{\,s},\ggen_D)$, $X = \Tag_\gamma^{s}$, and $\com_\gamma =
\gen_3^{\gamma}\cdot\gen_4^{r_\gamma}$. The equalities for $\com$, $U$, $X$,
$\com_\gamma$, and the range of $\gamma$ hold by construction, and
$\pairing(U,\ggen_D)^{\utk'} = \pairing(\HE(\epoch)^{s\utk'},\ggen_D) =
\pairing(\ek^{\,s},\ggen_D) = \Pi$ gives the one for $\Pi$. For the tag
equality, $\Tag_\gamma = \pairing(\ek, \ggen_D^{1/(\utk''+\gamma)})$ gives
$\Tag_\gamma^{\utk''+\gamma} = \pairing(\ek,\ggen_D)$, which gives the
following:
%
  \begin{equation*}
%
    X^{\utk''}\cdot X^{\gamma} = \big(\Tag_\gamma^{\utk''+\gamma}\big)^{s} =
\pairing(\ek,\ggen_D)^{s} = \pairing(\ek^{\,s},\ggen_D) = \Pi.
%
  \end{equation*}

  \paragraph{Soundness.} The sigma protocol of~\Cref{prot:txid-tag-sigma}
extracts $(\utk', \utk'', \openinfo, \gamma, s, r_\gamma)$ from two accepting
transcripts sharing a first message. The epoch key is not extracted directly.
We instead fix $\ek := \HE(\epoch)^{\utk'}$, which gives the second equality of
$\rel_7$ by definition, and the first equality is the first verification
equation. What remains is the tag equality, together with showing that the
extracted $\gamma$ is the value the range argument speaks about.

  From $U = \HE(\epoch)^{s}$ and $\Pi = \pairing(U,\ggen_D)^{\utk'}$ we get $\Pi
= \pairing(\HE(\epoch)^{s\cdot\utk'},\ggen_D) = \pairing(\ek^{\,s},\ggen_D)$.
With $X = \Tag_\gamma^{s}$, the tag equality then reads as follows:
%
  \begin{equation*}
%
    \pairing(\ek^{\,s}, \ggen_D) = \big(\Tag_\gamma^{s}\big)^{\utk''+\gamma},
\qquad \text{that is} \qquad \pairing(\ek, \ggen_D)^{s} =
\big(\Tag_\gamma^{\utk''+\gamma}\big)^{s}.
%
  \end{equation*}
%
  The verifier rejects unless $U \neq 1_{\group{1}}$, which forces $s \neq 0$,
as $\HE(\epoch)$ generates $\group{1}$ with overwhelming probability over the
random oracle. Cancelling $s$ leaves the following:
%
  \begin{equation}\label{eq:rel7prime-core}
%
    \pairing(\ek, \ggen_D) = \Tag_\gamma^{\utk''+\gamma}.
%
  \end{equation}
%
  The verifier also rejects unless $\Pi \neq 1_{\group{T}}$. As $\Pi =
\pairing(\ek^{\,s},\ggen_D)$ with $s \neq 0$, this forces $\ek \neq
1_{\group{1}}$ and $\pairing(\ek,\ggen_D) \neq 1_{\group{T}}$.
By~\Cref{eq:rel7prime-core} the right hand side is then not the identity either,
so $\utk'' + \gamma \neq 0 \bmod p$ and we can invert it as follows:
%
  \begin{equation*}
%
    \Tag_\gamma = \pairing(\ek,\ggen_D)^{1/(\utk''+\gamma)} =
\pairing\big(\ek,\, \ggen_D^{1/(\utk''+\gamma)}\big).
%
  \end{equation*}
%
  This is the tag equality of $\rel_7$. Both non-degeneracy checks have now been
used, and neither is redundant: the first ruled out $s = 0$ and the second ruled
out $\utk'' = -\gamma$.

  Finally, the fifth verification equation is $X^{\theta_2}\cdot
X^{\theta_\gamma} = F\cdot\Pi^{c}$. As $X^{\theta_2}\cdot X^{\theta_\gamma} =
X^{\theta_2+\theta_\gamma}$, this
equation on its own binds only the sum $\utk'' + \gamma$, which is all
that~\Cref{eq:rel7prime-core} needs. The two values are separated elsewhere. The
response $\theta_2$ is shared with the first verification equation, where the
binding property of the Pedersen commitment $\com$ pins $\utk''$ to the value
certified for the user, and given $\utk''$ the sum then pins $\gamma$. The
response $\theta_\gamma$ is in turn shared with the sixth verification equation,
$\gen_3^{\theta_\gamma}\cdot\gen_4^{\theta_r} = A_\gamma \cdot \com_\gamma^{c}$,
which opens $\com_\gamma$ to the extracted $\gamma$. As the range argument runs
on that same $\com_\gamma$, the value shown to lie in $[0,\Max)$ is the same
$\gamma$ that appears here. A prover who shifts $\utk''$ to $\utk'' + \delta$
and $\gamma$ to $\gamma - \delta$ still satisfies the fifth equation, but breaks
the first.

  \paragraph{Zero knowledge.} We construct a simulator $\simu$ that, given only
$\inst_7$ and without $\witness_7$, produces a transcript $(U, \Pi, X,
\com_\gamma, \algoproof_7')$ computationally indistinguishable from an honestly
generated one.

  \begin{description}

    \item[Simulating $(U, \Pi, X, \com_\gamma)$.] $\simu$ samples $s, r_\gamma
\sample \Zp$ and sets $U = \HE(\epoch)^{s}$ and $X = \Tag_\gamma^{s}$.  It can
compute both, as $\HE(\epoch)$ and $\Tag_\gamma$ are public. It sets
$\com_\gamma = \gen_3^{\gamma^\star}\cdot\gen_4^{r_\gamma}$ for an arbitrary in
range $\gamma^\star$. It then samples $\Pi \sample \group{T}$ uniformly, as it
cannot compute $\pairing(\ek^{\,s},\ggen_D)$.

      The only deviation from the real distribution is therefore $\Pi$. Writing
$\hat{\gen} = \pairing(\HE(\epoch),\ggen_D)$ and $\hat{\ek} =
\pairing(\ek,\ggen_D)$, both of which a distinguisher can compute, we compare
the two tuples below:
      %
      \begin{equation*}
      %
	\big(\hat{\gen},\, \hat{\ek},\, \hat{\gen}^{s},\, \hat{\ek}^{s}\big)
\quad \text{versus} \quad \big(\hat{\gen},\, \hat{\ek},\, \hat{\gen}^{s},\,
\Pi\big).
      %
      \end{equation*}
      %
      This is exactly a DDH instance in $\group{T}$, so hardness of DDH in
$\group{T}$ makes the change undetectable, and it does so whether or not the
distinguisher holds $\ek$. An authorised auditor does hold it, so this matters.
This is where withholding $\ek^{\,s}$ pays. Were the $\group{1}$ element
published, the simulator would have to sample it in $\group{1}$, and the
argument would need DDH in $\group{1}$ in the presence of $\ek$. That is
strictly stronger, since $\Pi$ is recoverable from $\ek^{\,s}$ by one pairing
but not the other way round. $X$ needs no separate assumption at all, being a
function of public data and of the simulator's own $s$.

    \item[Simulating the sigma protocol.] With $(U, \Pi, X, \com_\gamma)$ fixed,
$\simu$ runs the standard honest verifier simulator. It samples the challenge
$c$ and the responses $\theta_1, \theta_2, \theta_h, \theta_s, \theta_\gamma,
\theta_r \sample \Zp$ uniformly, and solves the verification equations for the
first message:
  %
      \begin{gather*}
  %
	A = \gen_2^{\theta_1}\cdot\gen_3^{\theta_2}\cdot\gen_4^{\theta_h}\cdot
\com^{-c}, \qquad B = \HE(\epoch)^{\theta_s}\cdot U^{-c}, \qquad M =
\pairing(U,\ggen_D)^{\theta_1}\cdot \Pi^{-c}, \\
    %
	C = \Tag_\gamma^{\theta_s}\cdot X^{-c}, \qquad F =
X^{\theta_2+\theta_\gamma}\cdot\Pi^{-c}, \qquad A_\gamma =
\gen_3^{\theta_\gamma}\cdot\gen_4^{\theta_r}\cdot \com_\gamma^{-c}.
    %
      \end{gather*}
  %
      By construction $\algoproof_7'$ satisfies every verification equation, and
since $c$ and the responses are uniform and independent, the simulated
transcript is distributed identically to a real one.

      The range argument for $0 \leq \gamma < \Max$ is simulated by the zero
knowledge simulator of the underlying Bulletproofs protocol~\cite{SP:BBBPWM18},
run on the same $\com_\gamma$ with an arbitrary committed value. The hiding
property of $\com_\gamma$, which holds as $r_\gamma$ is uniform and $\gen_4$ is
independent of $\gen_3$, means no distinguisher can tell which value was used.

  \end{description}

  \paragraph{The non-interactive variant.} For soundness, instead of rewinding
the prover we rely on the forking lemma to extract $\witness_7'$, and the
argument then follows exactly as above. For zero knowledge, the challenge $c$ is
a hash modelled as a programmable random oracle: $\simu$ samples $c$, computes
$(A, B, M, C, F, A_\gamma)$ exactly as above, and then programs $\hash(\inst_7',
A, B, M, C, F, A_\gamma) = c$. Since $\hash$ is a random oracle this is possible
and indistinguishable from normal use of the hash function, so the
non-interactive version is simulatable without rewinding.

\end{proof}

\subsection{Equivalence of $\rel_6$ and $\rel_6'$}\label{sec:rel6-equiv}

Write $b = {\gen_0}\cdot\gen_1^{\sk}\cdot\gen_2^{\utk'}\cdot
\gen_3^{\utk''}\cdot\gen_4^r$ for the base of the credential, so that a valid
BBS+ signature has $E = b^{1/(z+e)}$, and recall from~\Cref{sec:tx-auth} that
the prover sets $E' = E^{r_1}$, $r_3 = 1/r_1$, $d = b^{r_1}\cdot\gen_4^{-r_2}$,
$\bar{E} = E'^{-e}\cdot d\cdot\gen_4^{r_2}$, and $r'' = r - r_2 r_3 - \rho$.

\begin{proof}\leavevmode

  \paragraph{Correctness.} Substituting $d$ into $\bar{E}$ gives $\bar{E} =
E'^{-e}\cdot b^{r_1}$, and since $b = E^{z+e}$ we have $b^{r_1} = E'^{z+e}$, so
$\bar{E} = E'^{z}$ and hence $\pairing(E',\ggen_Z) = \pairing(E'^{z},\ggen) =
\pairing(\bar{E},\ggen)$, which is the second clause. The third clause,
$\bar{E}/d = E'^{-e}\cdot\gen_4^{r_2}$, is the definition of $\bar{E}$
rearranged.

  For the fourth clause, $d^{r_3} = (b^{r_1}\cdot\gen_4^{-r_2})^{1/r_1} =
b\cdot\gen_4^{-r_2 r_3}$, so that the following holds:
%
  \begin{equation*}
%
    d^{r_3} = \gen_0\cdot\gen_1^{\sk}\cdot\gen_2^{\utk'}\cdot\gen_3^{\utk''}
\cdot\gen_4^{r - r_2 r_3}.
%
  \end{equation*}
%
  With $\com = \gen_2^{\utk'}\cdot\gen_3^{\utk''}\cdot\gen_4^{\rho}$, taking
inverses and collecting the $\gen_4$ exponent as $-(r - r_2 r_3) + r'' = -\rho$
gives $\gen_1^{\sk}\cdot d^{-r_3}\cdot\gen_4^{r''} =
\gen_0^{-1}\cdot\gen_2^{-\utk'}\cdot\gen_3^{-\utk''}\cdot\gen_4^{-\rho} =
(\gen_0\cdot\com)^{-1}$, which is the fourth clause. The first clause is the
opening of $\com$ itself.

  \paragraph{Soundness.} The verifier checks $E' \neq 1_{\group{1}}$ and
$\pairing(E',\ggen_Z) = \pairing(\bar{E},\ggen)$ directly on public data, so
$\bar{E} = E'^{z}$. The sigma protocol of~\Cref{prot:txid-tx-nizk} extracts
$(\utk', \utk'', \rho, e, r_2, \sk, r_3, r'')$ from two accepting transcripts
sharing a first message. From the third clause, $E'^{z} = \bar{E} = d\cdot
E'^{-e}\cdot\gen_4^{r_2}$, so $E'^{z+e} = d\cdot\gen_4^{r_2}$. Raising to $r_3$
and substituting the fourth clause recovers a base of the required form, so
$(E'^{r_3}, e, \cdot)$ is a BBS+ signature on the messages committed in $\com$,
together with $\gen_1^{\sk}$.  The requirement $E' \neq 1_{\group{1}}$ rules out
the degenerate transcript in which every pairing equality holds vacuously.

  \paragraph{Zero knowledge.} $E'$ is uniform in $\group{1}$ as $r_1 \sample
\Zp^*$ is fresh per transaction, and $d$ is uniform independently of $E'$ as
$r_2 \sample \Zp$ is fresh, so $(E', \bar{E}, d)$ reveals nothing about
$\signature$, and $\bar{E}$ is determined by $E'$ and the public key. The sigma
protocol is simulated exactly as for $\rel_7'$ above, by sampling $c$ and
$\vec{\pi}$ and solving each verification equation for $t_1$, $t_2$, and $t_3$.

\end{proof}

\subsection{Proof of \texorpdfstring{\Cref{thm:inst}}{the instantiation
theorem}}\label{sec:txid-inst-sec}

\Cref{thm:txid-bb} reduces security of the scheme to five conditions on its
building blocks. It remains to discharge each of them for the instantiation
of~\Cref{sec:txid-instantiation}, which the preceding results do as follows.

\begin{proof}

\begin{description}

  \item[(P1).] $\DS$ is the BBS+ signature scheme of~\Cref{prim:bbs}, which is
existentially unforgeable under the $q$-SDH assumption (\Cref{asm:qsdh-bbs}).
The
variant used signs group elements rather than scalars, so unforgeability
additionally requires that the signer only sign elements accompanied by a proof
of knowledge of their representation. $\algoproof_4$ in $\IDregister$ supplies
this.

  \item[(P2).] $\Com$ is the Pedersen commitment scheme of~\Cref{prim:pedersen},
which is perfectly hiding, and computationally binding under the discrete
logarithm assumption in $\group{1}$. Binding holds against every party,
including the registration authority, because Step 2 of $\IDgen$ derives
$(\gen_2,\gen_3,\gen_4)$ by hash-to-curve from a public seed rather than letting
any party sample them.

  \item[(P3).] Pseudorandomness of $\FE$ is~\Cref{lem:txid-fe-prf}, under DDH in
$\group{1}$ with $\HE$ modelled as a random oracle. Pseudorandomness of $\FT$
is~\Cref{lem:txid-ft-prf}, under $q$-mDBDHI with $q = \Max$. Both statements
hand the distinguisher the whole of $\PTK$, including $\{W_\gamma\}$, which is
what Game 4 and Game 5 need, since every corrupt auditor holds it. The two
lemmas compose without a hybrid over epochs, and the simulator of Game 4
reproduces the real distribution by reusing one $\PTK_\user$ per user.

  \item[(P4).] Four proof systems are in play. First, the sigma protocols
of~\Cref{prot:txid-tag-sigma,prot:txid-tx-nizk}, Fiat-Shamir transformed with
$\hash$ modelled as a random oracle. Then the Schnorr proof of representation
used for $\rel_4$, and the batched proof for the $\{W_\gamma\}$ clause of
$\rel_5$. Last, the Bulletproofs range argument~\cite{SP:BBBPWM18}, whose
soundness rests on discrete logarithm in $\group{1}$. That the rewritten
relations prove what the
original ones state is~\Cref{sec:rel7-equiv,sec:rel6-equiv}, and extraction is
by rewinding, which is sound in the standalone model
of~\Cref{sec:txid-definitions}.

  \item[(P5).] Assumed of the network, as in the black-box construction.

\end{description}

Combining these with~\Cref{thm:txid-bb} gives~\Cref{thm:inst}.

\end{proof}
